\documentclass{article}
\usepackage[letterpaper, margin=1in]{geometry}
\usepackage{setspace}
\usepackage{indentfirst}
\usepackage{hyperref}
\usepackage{graphicx}
\usepackage{amsmath, amssymb}
\usepackage{bbold}

\usepackage[style=numeric, sorting=none, giveninits=true]{biblatex}
\addbibresource{bibs_final.bib}

\hfuzz=50pt

\title{Photon mediated interactions in open quantum systems}
\author{Sreeram Vasudevan}
\date{}
\begin{document}
\maketitle
\newpage
\tableofcontents

\doublespacing

\section{Introduction}

Photon-mediated interactions provide a powerful mechanism for coupling spatially
separated quantum emitters, enabling studies of fundamental light-matter interactions
and quantum technologies. When multiple emitters couple to the same photonic
reservoir, the exchange of virtual and real photons mediates both coherent
interactions and collective dissipation, giving rise to cooperative effects \cite{svidzinsky_cooperative_2010}. These interactions can result in drastic changes of the optical
response, known as Dicke superradiance or subradiance when the
ensemble is confined to regions of subwavelength light spacing \cite{dicke_coherence_1954,gross_superradiance_1982}. These effects are not limited to certain experimental platforms, they
are instead collective effects that follow from the interactions between the
electromagnetic field and the atoms, which can be explained via Green's functions
formalisms. These effects appear in similar forms irregardless of the environment
through which the electromagnetic field propagates. The most convenient
environments in quantum optics are free space, waveguides, or optical cavities.

Here, I examine the theoretical foundations of photon-mediated interactions, the
platforms to realize them, and advances toward programmable interaction networks. I
also take a look at possible challenges and a future outlook towards controllable
networks of interactions.

\section{Theory of photon-mediated interactions}

Quantum emitters rarely exist in a closed system where all of the information is
self-contained. More often, they interact with their surroundings and transfer
information to the environment. This transfer allows both for the long-range couplings
between emitters, but also serves as a channel for loss of photons. A complete
description of quantum emitters then requires the framework of open quantum systems.
In this section, I develop the model for describing inter-emitter interactions, show
how the electromagnetic field induces coherent and dissipative interactions, and
discuss effective descriptions that simplify resulting dynamics.

\subsection{Light matter interactions}

The two-level atom serves as the basis for simple models when considering
inter-emitter interactions. Each emitter $i$ is taken to have two levels, with a ground
state $|g_i\rangle$ and an excited state $|e_i\rangle$, separated by a transition
frequency $\omega_i$. This atom is then coupled to the quantized electromagnetic field
via the dipole interaction. This approximation is justified for atomic systems because
suitable isolated transitions can be identified, minimizing the effect of neighboring
energy levels on interaction dynamics. The total Hamiltonian for one two-level system
coupled to a continuum of field modes consists of the field contribution, the atomic
contribution, and the interaction between the atom and the field:
\begin{equation}
  \hat{H} = \hat{H}_{\mathrm{field}} + \hat{H}_{\mathrm{atom}} + \hat{H}_{\mathrm{int}}
\end{equation}

For a free field with wave number ${\bf k}$ and polarization $\mu$, we use creation
and annihilation operators to dictate the number of photons in each mode
\cite{cohen-tannoudji_atom-photon_1992, buhmann_dispersion_2013}:
\begin{equation}
\hat{H}_{\mathrm{field}} = \sum_{{\bf k},\mu} \hbar \omega_{\bf k}\left(\hat{a}_{{\bf
k},\mu}^\dagger
\hat{a}_{{\bf k},\mu}+\frac{1}{2}\right)
\end{equation}

The atomic contribution is
\begin{equation}
  \hat{H}_{\mathrm{atom}} = \frac{1}{2}\sum_i \hbar\omega_i\hat{\sigma}_z^{(i)}
\end{equation}
where we have the Pauli matrix $\hat{\sigma}_z = |e\rangle\langle e|-|g\rangle\langle
g|$.
The interactions between the atoms and the field are described by the atomic dipole
and the electric field the dipole interacts with \cite{buhmann_dispersion_2013,knoll_qed_2000,dung_resonant_2002}:
\begin{equation}
  \hat{H}_{\mathrm{int}} = -\sum_i {\bf d}_j\cdot{\bf E}_j
\end{equation}

The dipole operator couples the ground and excited states of each atom ${\bf d}_j =
d_{eg}\hat{\sigma}_+^{(j)}+d_{eg}^*\hat{\sigma}_-^{(j)}$, where $d_{eg}$ is the dipole moment
corresponding to the atomic transition. We also expand the electric field operator as
${\bf E}_j \propto \sum_{k,\mu}
e^{-i\omega_kt}\hat{a}_{k,\mu}+e^{i\omega_kt}\hat{a}^\dagger_{k,\mu}$.
We then make the rotating wave approximation to reach the interaction Hamiltonian for a single atom interacting with the field:
\begin{equation}
  \hat{H}_{\mathrm{int}} = \frac{\hbar}{2} \left(ge^{-i\Delta t}\hat{a} \hat{\sigma}_+ +
  g^*e^{i\Delta t}\hat{a}^\dagger \hat{\sigma}_-\right)
\end{equation}
where $g$ is the coupling strength between the dipole and the bath mode, $\Delta$ is
the detuning between the field mode and the atomic transition frequency, $\hat{a}$ and
$\hat{a}^\dagger$ are the photonic annihilation and creation operators, and the Pauli
matrices are defined $\hat{\sigma}_+ = |e\rangle\langle g|$,
$\sigma_-=|g\rangle\langle e|$. The rotating wave approximation is made here to
average out the counter rotating terms in the Hamiltonian (e.g.
$\hat{a}\hat{\sigma}_-$ and $\hat{a}^\dagger\hat{\sigma}_+$). These terms accumulate phase
too quickly to affect dynamics, as the system which accumulates phase much slower only
sees the average of the oscillations of these quickly rotating terms. This
approximation is valid under the assumption that the field mode interacting with the
atomic dipole is near resonant to the atomic transition.

The derived Hamiltonian provides a description of the coupled emitter-field system.
However, solving for the dynamics of a full system becomes difficult due to the
continuum of electromagnetic field modes. Since our interest lies in the emission of
the emitters, it would be advantageous to eliminate out the photonic degrees of
freedom.

\subsection{Master equation}
We can use the master equation to describe the dynamics of the system \cite{dzsotjan_quantum_2010,kastel_suppression_2005,meystre_elements_2007}:
\begin{equation}
  \dot{\rho} = -\frac{i}{\hbar}\left[\hat{H}, \rho\right] + \sum_i \Gamma_i\left(\hat{L}_i\rho
    \hat{L}_i^\dagger -
  \frac{1}{2}\left\{\hat{L}_i^\dagger \hat{L}_i, \rho\right\}\right)
\end{equation}
where $\Gamma_i$ is the coupling into the loss modes described by the operators $L_i$.
When we set the coupling to the loss modes $\Gamma_i = 0$, we recover unitary dynamics
$\dot{\rho}=-\frac{i}{\hbar}[H, \rho]$. By invoking the master equation, we make two
key approximations, which we attribute a single name: the Born-Markov approximation.
The Born approximation assumes that the bath is much larger than the system and that
any effect of the system on the bath is negligible. This allows us to simply take the
tensor product of the system and the bath to get the total density matrix. The Markov
approximation is that of time scales: it assumes the bath correlation time is
much smaller than that of the system relaxation time. If this is the case,
then the bath effectively has no memory of any history of the system, and evolution
only cares about the present state.

The master equation provides a general description of emitter dynamics, but the
coherent and dissipative couplings depend explicitly on the electromagnetic
environment. A general approach is provided by the electromagnetic Green's tensor,
which describes the response of an arbitrary medium to a point source.

\subsection{Green's function formalism}
We reframe the electric field operator in terms of the electromagnetic Green's
tensor. In this way, the electric field is written in terms of the bosonic
creation and annihilation operators \cite{gruner_green-function_1996,
asenjo-garcia_atom-light_2017}:
\begin{equation}
  E(r,\omega)^+ = i\mu_0\omega^2\sqrt{\frac{\hbar\epsilon_0}{\pi}}\int
  dr'\sqrt{\epsilon_I(r',\omega)}G(r, r', \omega)\hat{f}(r',\omega)
\end{equation}

The total electric field $E(r,\omega) = E^+ + E^-$, where $E^{+\dagger} = E^-$. Here, we define the function $\epsilon = \epsilon_R + \epsilon_I$ as the complex
permittivity. The operators $\hat{\mathbf{f}}^\dagger$ and $\hat{\mathbf{f}}$ are the
bosonic creation and annihilation operators corresponding to the electric field.
Unlike the operators used in the free-space quantization, these operators incorporate
both the electromagnetic field and reservoir degrees of freedom, ensuring the field
satisfies the proper bosonic commutation relationships while also accurately
describing the dissipative environment.
The classical Green's
function $G(r, r', \omega)$ describes the propagation of the electromagnetic field
through the medium. In this treatment of the electric
field, we assume that the only interactions between dipoles are those mediated by this
medium-assisted electromagnetic field.

We can then redefine the total field Hamiltonian in terms of these bosonic operators
\cite{huttner_quantization_1992,knoll_qed_2000}:
\begin{equation}
  H_{\mathrm{field}} = \int d{\bf r}\int d\omega\, \hbar\omega
  \hat{\mathbf{f}}^\dagger(\mathbf{r},\omega)\hat{\mathbf{f}}(\mathbf{r}, \omega)
\end{equation}


The emitters interact with via processes mediated by the photonic reservoir. For
coherent processes, the photon is emitted into the reservoir by one emitter and is
then absorbed by another emitter. This exchange interaction is the transfer of a
virtual photon through the bath. Emitters can also dissipatively couple with one
another, in a process where both atoms emit photons which constructively or
destructively interfere, leading to phenomena such as subradiance and superradiance.

Under this framework, we again trace out photonic degrees of freedom, apply the
Born-Markov approximation, and retrieve the effective Hamiltonian
\cite{dung_resonant_2002, chang_cavity_2012,chang_trapping_2014}:
\begin{equation}
  H = -\hbar \Delta_A \sum_{i=1}^N \sigma_{ee}^i - \hbar\sum_{i,j=1}^N
  J^{ij}\sigma_+^i\sigma_-^j - \sum_{i=1}^N[\mathbf{d} \cdot
  \hat{\mathbf{E}}^-(\mathbf{r}_i)\sigma_-^i + \mathbf{d}^*\cdot
  \hat{\mathbf{E}}^+(\mathbf{r}_i)\sigma_+^i]
\end{equation}

We also can retrieve the Linbladian term
\begin{equation}
  L[\rho] = \sum_{i.j=1}^N \Gamma^{ij}\left( \sigma_-^i\rho\sigma_+^j - \frac{1}{2}
  \left\{\sigma_+^i\sigma_-^j, \rho\right\} \right)
\end{equation}

From these two equations, we can retrieve two important terms: the spin exchange rate
$J^{ij}$ and the collective decay rates $\Gamma^{ij}$. They can also be written in
terms of the Green's functions as:
\begin{align}
  J^{ij} &= (\mu_0\omega^2/\hbar) \mathbf{d}^*\cdot \mathrm{Re}\mathbf{G(r}_i, {\bf
  r}_j, \omega)\cdot {\bf d}\\
  \Gamma^{ij} &= (2\mu_0\omega^2/\hbar) {\bf d}^* \cdot \mathrm{Im}\mathbf{G(r}_i,
  {\bf r}_j, \omega)\cdot {\bf d}
\end{align}

Here we see how the coherent exchange processes and dissipative interactions between
emitters are related to both the dipole elements of the transition of the atom you are
using, as well as the Green's function provided by the medium through which the
electromagnetic field propagates.



\section{Experimental platforms}
The formalism derived in Section 2 claims that every platform that couples emitters to
a shared photonic environment is described by the same set of equations, differing in
the form of the Green's functions, provided that the assumptions we made still hold.
This section delves further into that claim and showcases some particular experimental
realizations and applications of photon-mediated interactions. The photonic platforms
of interest are free space, waveguides, and optical cavities.

\subsection{Free space arrays}

In free space, emitters are held at fixed positions with no engineered photonic
structure between them. The relevant $\mathbf{G}$ is the free space dyadic Green's
function, which is of the form \cite{novotny_principles_2012}
\begin{equation}
  \mathbf{G(r}) = \frac{k^2e^{ikr}}{4\pi\epsilon_0r}\left[ \left(1 + \frac{i}{kr} -
    \frac{1}{(kr)^2} \right) \mathbb{1}+ \left(-1-\frac{3}{kr} + \frac{3}{(kr)^2}\right)
\hat{\mathbf{r}}\otimes\hat{\mathbf{r}} \right]
\end{equation}
which has an oscillatory component in $r$ as well as terms weighted by $1/kr$,
$1/(kr)^2$, and $1/(kr)^3$. The near field terms dominate at short separation
distances and fall off steeply, so free-space coupling is short range by default. In
this regime, we are able to see interesting collective phenomena such as superradiance
and subradiance. 

Ordered arrays of atoms in an optical lattice can be used to demonstrate some
interesting effects. For example, a group was able to use subradiant modes that appear in
sub-wavelength spaced Rubidium atoms in a 2-d lattice to create a single atom thick
mirror for photons \cite{rui_subradiant_2020}. Figure \ref{fig:reflectance} shows how the reflectance of the 2d
array increases as the probe laser gets closer to resonance. As previously mentioned
in the theory section, the collective effects are much stronger when the atoms are
driven on or near resonance. This paper further states that directly driven subradiant
modes require a spatially ordered 2d array, which comes from the coupled dipole
theory.

\begin{figure}[!htbp]
  \centering
  \includegraphics[width=0.4\linewidth]{aaaaa}
  \caption{Transmittance and reflectance for the excitation of a 2d array. The
  response is subradiant, and the transmission drops and the reflectance increases as
the probe gets closer to resonance. The 2d array, when driven closer to resonance
starts to exhibit collective phenomena. Reproduced from \cite{rui_subradiant_2020}.}
  \label{fig:reflectance}
\end{figure}

We also explore another paper by Ruttley, et. al. where molecules in optical tweezers
were used in entanglement experiments \cite{ruttley_long-lived_2025}. Ultracold polar molecules were trapped in
arrays of optical tweezers, and then long-range interactions were engineered. In the
previous example, we only considered sub-wavelength spacing, but here the molecules
are spaced at ranges where generally only hertz-scale interactions dominate. By using
microwave photons, they were able to couple two molecules trapped in magic-wavelength
tweezers,
and demonstrated a high level of entanglement fidelity of about 97.6\%.

\subsection{Waveguide}
In a waveguide, emitters couple to one or more guided modes of an effectively
one-dimensional structure. These structures can range from optical fibers to photonic
crystal waveguides to superconducting transmission lines. The guided mode of a
waveguide propagates without transverse spreading, and thus the guided mode
contribution of the Green's function is purely oscillatory \cite{barton_elements_1989,tai_dyadic_1971}:
\begin{equation}
  {\bf G} \propto e^{ik|z_i - z_j|}
\end{equation}
where $z_i$ denotes the position of emitter $i$.

The non-decaying Green's function is what makes waveguides unique. The coherent
and dissipative couplings between emitters remain non-negligible at long distances and
is limited by propagation loss and disorder. The relative phase can be tuned by way of
emitter spacing, and thus we are able to easily choose whether coupling is
superradiant or subradiant.

Superradiance can be found for atoms trapped along a photonic crystal waveguide (PCW).
Goban, et. al. used an alligator photonic crystal waveguide to trap atoms and
retrieve the collective dynamics that emerge \cite{goban_superradiance_2015}. These interactions provide a proof of
concept for the idea of trapping cold atoms near PCWs in order to achieve
superradiance. PCWs can then be used as a platform for manipulating collective
effects, with applications in quantum matter and quantum simulation.

We can take a further look at long range interactions from Tiranov, et. al.,
where they study super- and subradiant dynamics between distant quantum emitters
\cite{tiranov_collective_2023}. In
this experiment, quantum dots are placed at distances of between 4 and 5 times the
wavelength within the medium and the photonic crystal waveguide. One of the quantum
dots is then excited, and the collective emission is then measured. Off resonance, the
quantum dots emit as if they were independent, but as the drive gets closer to
resonance, the collective emission shows the effects of both super- and subradiance.
Super- and subradiance in this experiment originate from the constructive and
destructive emission of the field emitted into the waveguide, and is seen by the rate
of decay of the photon counts at the detectors.

As mentioned previously, these effects are not platform specific, but rather a
hallmark of photon mediated interactions.

\begin{figure}[!htbp]
  \centering
  \includegraphics[width=0.4\linewidth]{abcde}
  \caption{Measured emission for a pair of emitters spaced by 1.25(3)$\mu$m both on
  and off resonance. We see the collective effects of subradiant emission as we move
closer to resonance. Reproduced from \cite{tiranov_collective_2023}}
  \label{fig:pcw}
\end{figure}

Superconducting qubits, as well as the superconducting transmission lines are a great
choice of waveguide in order to do this type of QED. Superconducting qubits have
extremely high coupling coefficients to the transmission lines, and the lines behave
like waveguides do, and allow for long range coupling \cite{mirhosseini_cavity_2019,astafiev_resonance_2010}. This is particularly useful for
quantum information as we can create states that have long lifetimes, and we can also
read them out. For example, Zanner, et. al. have been able to show coherent control
over multi-qubit dark states in this type of waveguide QED \cite{zanner_coherent_2022}. In these systems, we are
more interested in the states that cause subradiant emission, known as dark states.
Dark states are very useful in quantum information as their long lifetimes are useful for
quantum simulation and quantum computation \cite{lidar_decoherence-free_1998,paulisch_universal_2016}. They were able to create a four-transmon dark state which had a relaxation
rate 160 times lower than that of the single qubit rate.
\subsection{Optical cavities}
In cavities, emitters couple mainly to a single cavity mode, set by the structure of
the cavity. The Green's function is that of a standing wave, and is
sinusoidal \cite{lodahl_interfacing_2015}. This allows all atoms to interact since they all see the same cavity mode.
We can engineer the coupling strength as well as the ratio of coherent to dissipative
couplings. For example, trapping the atoms at the peaks of the standing wave leads to
a strong coupling to the cavity.

An example of using cavities to demonstrate collective effects was the
Thompson group, who were able to create a superradiant laser with less than one intracavity
photon \cite{bohnet_steady-state_2012}. By trapping rubidium atoms in an optical cavity, they are able to create a
narrow linewidth laser. The collective atomic dipoles store the coherence, and we are
still able to get stimulated emission even with less than one photon in the cavity.

\section{Engineering interaction networks}
One interesting direction to look towards in the regime of photon mediated
interactions are engineering arbitrary interaction networks. These systems allow for
the study of more complicated structures, such as spin glass phases \cite{gopalakrishnan_frustration_2011,strack_dicke_2011} or
topologically ordered states \cite{hung_quantum_2016}. With most of the
methods mentioned in the previous sections, interactions are all-to-all. Any pair of
emitters are able to couple due to them all sharing a photonic reservoir. If we wanted
to create a network with select interactions turned off, the current model does not
suffice. Finding a solution for this question is a direction some
groups have recently been looking towards, where we now intend to create networks of
interactions that are useful for applications such as quantum simulation or quantum
sensing.

\subsection{Effective descriptions}
The formalism we have thus far used detailed the process by which we can describe
dynamics in an open quantum system with simple two-level systems coupled to some
electromagnetic bath. 
Now, we wish to expand our toolbox to describe systems which can
have multiple decay channels. We can use systems with multiple decay channels to
facilitate selective spin-exchange interactions \cite{periwal_programmable_2021}.
If we have an atom which is not strictly a two level system, it can
have multiple excited states or multiple ground states, which increases the possible
number of channels a photon can travel through and be emitted by. Reiter and Sorensen
(2012) \cite{reiter_effective_2012} provide a method by which we can create an effective
master equation for systems with multiple decay processes.

The proposed method allows for the adiabatic elimination of the excited manifold,
lowering the complexity of the system by creating an effective master equation of the
ground states. We are able to adiabatically eliminate out the excited states because
when the excited state is driven far off resonance or has very short decay times, the
average population that lives in the excited state is almost zero. When applied to our
system of emitters, this leads to a system in which we are left with jump operators
acting solely on the ground manifold. This also leaves us with an effective spin model
with interactions of the form $J^{ij}$ and $\Gamma^{ij}$, as in the previous section.

As opposed to traditional adiabatic elimination, which is generally done at the level
of the equations of motion, this method works on the operators and gives the effective
operators without having to revert to the equations of motion.

The effective Hamiltonian, in the effective description can be written
\begin{equation}
  H_{\mathrm{eff}} = H_g - \frac{1}{2}\left[ \hat{V}_- \sum_{f,l}
  \left(H_{\mathrm{NH}}^{(f,l)}\right)^{-1}\hat{V}_+^{(f)} + \mathrm{H.c.} \right]
\end{equation}
where $H_g$ is the Hamiltonian that describes the ground state manifold of the system,
$H_e$ describes the excited manifold, and $V_+$ and $V_-$ are the couplings between
the ground and excited subspaces, where $V_+$ takes you from the ground subspace to
the excited subspace, and $V_-$ the opposite. $H_{\mathrm{NH}}$ is the non-Hermitian
Hamiltonian which is constructed with respect to each ground state $f$ and each
coupled field mode $l$.

We also derive the effective jump operator in a similar manner.
\begin{equation}
  L_{eff}^k = L_k \sum_{f,l} \left(H_{\mathrm{NH}^{(f,l)}}^{-1}\right)V_+^{(f)}
\end{equation}

This formalism allows us to go further beyond the emitter-bath picture derived in the
previous section by allowing for further complications in the emitter itself.

\subsection{Arbitrary interaction networks}
The group of Schleier-Smith studied an ensemble of atoms in a cavity in order to
engineer programmable interaction networks \cite{periwal_programmable_2021}. By introducing a magnetic field gradient along the
cavity axis, they introduce position dependent Zeeman splittings, making spin-exchange
processes between atoms at different locations off-resonant.
Then, by modulating the drive frequency of the cavity, they were able to create energy
sidebands that compensated for the energy mismatches. This allows them to selectively
restore spin exchange correlations between chosen pairs of atoms, and engineer
interaction graphs beyond the bare all-to-all cavity mediated coupling.

\begin{figure}[!htbp]
  \centering
  \includegraphics[width=0.7\linewidth]{mss}
  \caption{Correlations for different combinations of magnetic field gradient.
  Interactions stop becoming all-to-all, and we can keep only local interactions, or
keep interactions that exist only at specific distances. Reproduced from
\cite{periwal_programmable_2021}}
  \label{fig:mss}
\end{figure}

This particular method works for creating interesting coherent networks, but this does
not allow for the creation of dissipative networks with similar structures. Since the
interesting couplings happen due to spin exchange processes, this cannot lead to a
dissipative model.

\section{Challenges and Outlook}

The theoretical picture developed in Section 2 and the experimental platforms
described in Section 3 show that coherent and dissipative photon-mediated interactions
are both derived from the same initial formalism, and the platforms are able to
exhibit the collective effects that arise from ensembles of coupled emitters. However,
using this theoretical picture to generate a programmable network of interactions is a
difficult problem.

The shared-bath picture that allows for photon-mediated interactions happens to also
be the big drawback when it comes to trying to program different networks. Since every
emitter couples to the same reservoir modes, the default network is all-to-all, with
the exact strengths and forms of the couplings set by the environment. The magnetic
field gradient technique used by the Schlier-Smith group demonstrates one possible
route to selectivity. Introducing site-dependent energy shifts can bring only chosen
pairs into resonance, but this approach is inherently coherent. Reproducing the same
programmability for purely dissipative networks is still an open question.

The issue with dissipative interaction networks in this current formalism is that all
of the emitters and their emissions interact with all of the same reservoir modes. If
we were able to modify pairs of emitters such that only their emissions would
interact, we could then engineer arbitrary networks. One answer
would be to have separate, non-interacting modes which only couple to those emitters
we would want to interact \cite{nair_topological_2023}. However, in practice, engineering large amounts
of non-interacting modes to generate nearest-neighbor couplings becomes difficult. We
can instead modify emissions off of each atom by introducing a dressing laser, leading
to only emissions of the same frequency interacting
\cite{vasudevan_engineering_nodate}. This process is 
dissipative interaction networks.


\addcontentsline{toc}{section}{References}
% \bibliographystyle{abbrv}
% \bibliography{bibs_final}
\printbibliography
\end{document}
